% =====================================================================
% Case study II: Dedekind delegation (Sec 7) - sources:
%   * CourseChain/Dedekind.lean (header design rules; per-theorem
%     mathlib references: IsPrincipalIdealRing.isDedekindDomain,
%     Ideal.dvd_iff_le, Ideal.uniqueFactorizationMonoid,
%     Ideal.prod_normalizedFactors_eq_self, Ideal.prime_iff_isPrime,
%     Ideal.absNorm, Ideal.absNorm_span_singleton,
%     IsDedekindDomain.primesOver_finite; delegation modules
%     Mathlib.RingTheory.DedekindDomain.{Basic,Ideal.Basic,Ideal.Lemmas};
%     textbook anchor Neukirch GTM 322 ch. I & III.1)
%   * framework/proof_status.md sec 3 (THEOREM registration +
%     three-part rationale; FLT/Langlands THEOREM* rows)
%   * framework/journal_submission_d1_review.md sec 2.4-4 ("zero
%     self-built mathematics, validating the mathlib attribution
%     mechanism"), 2.4-7 (attribution table normalization)
%   * git log v7.95/v7.96 (course-chain segments Zp/Qp, Dirichlet)
% Number policy: all numbers document-anchored; no new calibers.
% =====================================================================

\subsection{The course chain and the delegation pattern}
\label{sec:dedekind-chain}
The second case study is deliberately the opposite of the first.
The Chern--Simons chain (Sec.~\ref{sec:cschain}) is a physics claim
whose deepest provable layer had to be \emph{found} inside a
dependency chain of placeholders; the Dedekind module is a textbook
chapter of commutative algebra whose proofs were never in doubt and
whose only governance question is \emph{attribution}: when every
theorem in a file is discharged by a named mathlib lemma, what may
the library claim about itself?

\emph{CourseChain/Dedekind.lean} is the second segment of a course
chain (first segment: $\mathbb{Z}_p/\mathbb{Q}_p$, v7.95; this
segment: Dirichlet theory background plus Dedekind domains, v7.96)
that formalizes the algebraic-number-theory curriculum in Lean~4.
It states the foundational theorems of Dedekind domain theory ---
a PID is a Dedekind domain; $\mathbb{Z}$ is a Dedekind domain;
ideal divisibility is reverse containment; unique factorization of
ideals; prime ideals and residue norms; finiteness of primes over
a given ideal --- and proves each by \emph{direct delegation}:
\texttt{exact}/\texttt{:=} to a named mathlib lemma or instance
whose signature was verified against the pinned mathlib revision
(\texttt{8a178386}) in the module header. The design rules are
explicit: no unfinished proofs, no new postulates, no trivial
placeholders, and \emph{only} theorems whose mathlib counterpart
could be located and signature-verified --- ``no guessing of lemma
names.'' The textbook anchor (Neukirch, \emph{Algebraic Number
Theory}, Chapters~I and~III.1) is cited in the header so that the
curriculum intent is checkable against the literature.

Table~\ref{tab:dedekind-attr} is the attribution table the module
maintains, normalized to paper form (per-lemma mathlib anchors
recorded in the source; the flagship instance chain shown).

\begin{table}[t]
\caption{Dedekind delegation table (CourseChain/Dedekind.lean):
every theorem discharges to a named mathlib entity; the flagship
row carries the full instance chain. Zero \texttt{sorry}, zero
\texttt{axiom}, zero \texttt{trivial} in the module.}
\label{tab:dedekind-attr}
\begin{center}
\small
\begin{tabular}{@{}p{4.6cm}p{6.0cm}p{3.2cm}@{}}
\toprule
Course-chain theorem & mathlib delegate & Delegate kind \\
\midrule
\texttt{pid\_isDedekindDomain} &
\texttt{IsPrincipalIdealRing.\allowbreak isDedekindDomain} & instance \\
\texttt{int\_isDedekindDomain} & (as above, via
\texttt{inferInstance}) & instance \\
\texttt{ideal\_dvd\_iff\_le} & \texttt{Ideal.dvd\_iff\_le} &
lemma \\
\texttt{ideal\_\allowbreak uniqueFactorizationMonoid} &
\texttt{Ideal.\allowbreak uniqueFactorizationMonoid}
(\texttt{Ideal.Basic} L432--447) & instance \\
\texttt{ideal\_prod\_\allowbreak normalizedFactors} &
\texttt{Ideal.prod\_\allowbreak normalizedFactors\_\allowbreak eq\_self} & lemma \\
\texttt{ideal\_prime\_iff} & \texttt{Ideal.prime\_iff\_isPrime}
& iff-lemma \\
\texttt{ideal\_absNorm} \emph{et al.} &
\texttt{Ideal.absNorm}, \texttt{Ideal.absNorm\_\allowbreak span\_singleton}
& def/lemma \\
\texttt{primesOver\_finite} &
\texttt{IsDedekindDomain.\allowbreak primesOver\_finite} & lemma \\
\bottomrule
\end{tabular}
\end{center}
\end{table}

\subsection{Zero new mathematics, and why it is registered anyway}
\label{sec:dedekind-zero}
The internal peer review of this module was blunt
(journal\_submission\_d1\_review, sec 2.1~\cite{sylvad1review2026}): ``all theorems are
delegation wrappers; zero new mathematics.'' Under a
mathematics-first venue that is a rejection. Under the governance
lens of this paper it is the cleanest possible demonstration of
what the claim tiers are \emph{for}:

\begin{itemize}
\item The registry entry is THEOREM, not THEOREM*: the proof
(compiled delegation to a proved mathlib instance chain) is
machine-checked within the corpus, complete, and carries zero
local obligations --- \texttt{inferInstance} to
\texttt{Ideal.\allowbreak uniqueFactorizationMonoid} (whose own derivation in
mathlib runs through ideal invertibility and
\texttt{WfDvdMonoid} to \texttt{UniqueFactorizationMonoid}),
with the pinned-commit source lines recorded. The three-part
registration rationale (statement's role in the curriculum; proof
chain with named delegates; compile-status basis) is the same
format as T3's.
\item THEOREM* is reserved for the other delegation pattern:
mathematics proved \emph{elsewhere} and formalized
\emph{elsewhere}, which the corpus only cites --- the registry's
Fermat's Last Theorem row (THEOREM*, ``mathlib-provided, not in
this repository'') is the canonical example, with the repository
quoting mathlib's scale for it (\pnum{12000} lines of FLT-specific
development atop a \pnum{270000}-theorem prerequisite stack,
per the registration note).
\item What the module \emph{adds} is not mathematics but a
verified attribution mechanism: signature-level anchoring of every
course-chain statement to a pinned mathlib revision, a textbook
correspondence, and a compile-clean module boundary. When the
course chain later grows segments whose proofs are only partly in
mathlib, the delegation table is precisely the surface on which
``local work vs.\ borrowed work'' will be visible --- the same
surface the five-class taxonomy provides for axioms.
\end{itemize}

\subsection{T3 vs.\ Dedekind: two ends of one scale}
\label{sec:dedekind-contrast}
Set side by side, the two case studies calibrate the whole
program. Both registrations are THEOREM; both rest on
\texttt{\#print axioms}-clean compilations against the same pinned
mathlib revision; both use \emph{named} delegation
(\texttt{Int.cast\_sum}/\texttt{Int.cast\_mul} for T3,
\texttt{inferInstance}/\texttt{exact} for Dedekind). They differ
in exactly one dimension: where the local contribution sits.
Dedekind's contribution is the attribution layer itself (statements
curated, signatures verified, zero local proof content); T3's
contribution is a small but real assembly step --- the integrality
of a stratified sum is not a mathlib lemma, and discharging it
required combining three named lemmas in a proof whose shape was
not copied from anywhere. The scale between ``pure delegation''
and ``assembly'' is the honest axis along which the corpus's
non-placeholder content must be presented; pretending either end
is ``a new theorem'' is exactly the failure mode F1 that the
governance program exists to prevent. The course chain continues
(v7.95: $\mathbb{Z}_p/\mathbb{Q}_p$; v7.96: Dirichlet background +
Dedekind), and each future segment will be graded on this axis at
registration time.
